\documentclass[11pt]{article}
\usepackage{mathblatt}
% Präambel des Setzers (werkzeuge/setzer-praeambel.tex, wortgleich); dazu \gteil (Muster T2U)
\makeatletter
\setlength{\mbpfbe}{0mm}% keine Punkte (pruefung.md § 5: nur beim Prüfstein)
\newlength{\gEin}\setlength{\gEin}{0pt}
\let\gSchrift\relax
\renewcommand{\pfkreuzzeile}[1]{\par\noindent\hangindent3.2em\hangafter1\hspace*{1.6em}$\square$\ #1\par}
\newcommand{\gkreuz}[1]{$\square$\ #1\hspace{2em}}
\newcommand{\gkreuzzeile}[2]{\par\noindent\hangindent2.2em\hangafter1\mbox{#1\,$\square$}\ #2\par}% Buchstabe, Option: eine je Zeile (neben Skizze)
% links, was man ansieht, rechts, was man tut; Spaltengrenze überall an derselben Stelle (gemeinsam.md „Satz“)
\newsavebox{\gbox}\newlength{\gLinks}
% v0.10 (Lehrer 13.10.): keine starre Halbierung – die linke Spalte ist so breit wie die Figur (+ 6 mm), die
% Antwort steht direkt daneben; ist die Figur breiter als 0,62 Zeile, steht die Antwort darunter
\newcommand{\gzwei}[2]{\par\smallskip\noindent\sbox{\gbox}{#1}%
  \setlength{\gLinks}{\dimexpr\wd\gbox+6mm\relax}%
  \ifdim\gLinks>0.62\linewidth
    \usebox{\gbox}\par\smallskip\noindent\begin{minipage}{\linewidth}#2\end{minipage}%
  \else
    \begin{minipage}[t]{\gLinks}\vspace{0pt}\usebox{\gbox}\end{minipage}%
    \begin{minipage}[t]{\dimexpr\linewidth-\gLinks\relax}\vspace{0pt}\raggedright #2\end{minipage}%
  \fi\par}
% v0.9 (12.10., pruefung.md § 5 „Marke mit B, Bedingungen rechts“): drittes Argument von pfaufg = Bedingungen
% der Prüfung (v0.10: \pfohnetr = TR durchgestrichen, „FOR“), klein grau im rechten Rand auf der Grundlinie der Nummer, Zeilen mit \\
\newcommand{\gbedrand}[2]{\ifblank{#2}{}{\rlap{\hspace*{\dimexpr#1+2mm\relax}%
  \parbox[t]{15mm}{\raggedright\scriptsize\leavevmode\color{mbgrau}\strut#2}}}}
% Teil eines Bündels (a, b …): Marke links im Rand, Bedingungen rechts im Rand, auf der Zeile des Teils
\newcommand{\gteilrand}[2]{\llap{\mbpfrandmarke{#1}\hspace*{\mbpfnr}}\gbedrand{\linewidth}{#2}}
% Teil eines Bündels: Buchstabe in eigener Spalte, Inhalt eingerückt (\gEin = Einzug, damit die Spaltengrenze von
% \gzwei an derselben Stelle bleibt)
\newenvironment{gbteil}[3]{\par\smallskip\noindent\gteilrand{#2}{#3}\makebox[\mbpfteil][l]{\textbf{#1}}%
  \begin{minipage}[t]{\dimexpr\linewidth-\mbpfteil\relax}\setlength{\gEin}{\mbpfteil}\raggedright\strut\ignorespaces}%
 {\end{minipage}\par}
% Grafik höchstens so breit wie die Spalte (Bündel nebeneinander, Koordinatensystem ohne Karo)
\newsavebox{\gpbox}
\newcommand{\gpassend}[1]{\sbox{\gpbox}{#1}\ifdim\wd\gpbox>\linewidth\typeout{GPASSEND: verkleinert \the\wd\gpbox\space auf \the\linewidth}\resizebox{\linewidth}{!}{\usebox{\gpbox}}\else\usebox{\gpbox}\fi}
% Skizzenfeld (v0.10, Lehrer 13.10.): Karo ab dem linken Rand über die Zeile, „Skizze“ klein grau oben links im Feld
% (wie \pfskizzenfeld der Vorlage); der rechte Rand bleibt den Bedingungen
\newcommand{\gskizze}[1]{\par\smallskip\noindent\begin{tikzpicture}
  \clip (0,0) rectangle ({\linewidth-1mm},{#1*10mm});
  \draw[black!16,line width=0.3pt,step=5mm] (0,0) grid ({\linewidth},{#1*10mm});
  \node[anchor=north west,inner sep=2pt,font=\scriptsize,text=mbgrau,fill=white] at (0,{#1*10mm}) {Skizze};
  \end{tikzpicture}\par}
\RenewDocumentEnvironment{pfaufg}{m m m}{%
  \begin{lrbox}{\mbaufgabebox}\begin{minipage}[t]{\linewidth}%
  \gSchrift\noindent\hspace*{\gEin}\mbpfrandmarke{#1}\gbedrand{\linewidth-\gEin}{#3}%
  \begin{minipage}[t]{\mbpfnr}\strut\textbf{#2}\end{minipage}%
  \begin{minipage}[t]{\dimexpr\linewidth-\gEin-\mbpfnr\relax}\raggedright\strut\ignorespaces}%
 {\end{minipage}%
  \end{minipage}\end{lrbox}%
  \setlength{\mbaufgabehoehe}{\dimexpr\ht\mbaufgabebox+\dp\mbaufgabebox\relax}%
  \ifdefempty{\gkopfnach}{\par\addvspace{5pt}\Needspace*{\mbaufgabehoehe}}{%
    % Stufenkopf, der mit dieser (ersten) Aufgabe zusammenbleibt: Platz für beide verlangen, dann den Kopf setzen
    % (12.10., Quadratische: ein Kopf stand allein unten, die Aufgabe mit Koordinatensystem auf der nächsten Seite)
    \par\Needspace*{\dimexpr\mbaufgabehoehe+40pt\relax}\gkopfnach\gdef\gkopfnach{}\par\addvspace{5pt}}%
  \noindent\usebox{\mbaufgabebox}\par\addvspace{5pt}}
% Stufenkopf vormerken: \gkopfvor{<Kopf>} – die nächste pfaufg setzt ihn, sobald beide zusammen Platz haben
\def\gkopfnach{}
\newcommand{\gkopfvor}[1]{\gdef\gkopfnach{#1}}
% Teilaufgabe mit Skizze links in gzwei: Buchstabe in eigener Spalte, darunter Text und Skizze
\newcommand{\gteil}[2]{\makebox[\mbpfteil][l]{\textbf{#1}}\parbox[t]{\dimexpr0.5\textwidth-\gEin-\mbpfnr-5mm-\mbpfteil\relax}{\raggedright #2}}
% Stufenkopf der Vorlage (mathblatt.sty 2026-10-12: Linie, Abstand, 13 pt, \pfjahre) ohne ihren Vorhalt von
% 14 Zeilen: der Setzer fordert vor dem Kopf 0,16 Seitenhöhe (Kopf + erste Aufgabe), mehr nicht
\setlength{\mbpfstufevorhalt}{0pt}
\makeatother
\newcommand{\karo}[1]{\par\smallskip\noindent\begin{tikzpicture}
  \clip (0,0) rectangle ({\linewidth-1mm},{#1*10mm});
  \draw[black!16,line width=0.3pt,step=5mm] (0,0) grid ({\linewidth},{#1*10mm});
  \end{tikzpicture}\par}
\newcommand{\kk}[1]{$\square$\,#1\hspace{0.9em}}
\newcommand{\linie}[1]{\,{\color{black!45}\rule[-1pt]{#1}{0.4pt}}\,}
\newcommand{\pkt}{\textperiodcentered{}}
% v0.7 (12.10.): Schülerfrage unter dem Stufenkopf; Ankreuzoption als Bild (Buchstabe und Kästchen darunter)
\newcommand{\gfrage}[1]{\par\nobreak\noindent\hspace*{\mbpfnr}{\small\itshape #1}\par\nobreak\addvspace{2pt}}
\newlength{\gbildbreite}\setlength{\gbildbreite}{0.32\linewidth}% Setzer setzt sie je Reihe (vier Bilder: 0,24)
\newcommand{\gbildoption}[3]{\begin{minipage}[t][#3][b]{\gbildbreite}\centering #2\par\smallskip\mbox{#1\,$\square$}\end{minipage}}
% Lineare (12.10.): Bild mit Buchstabe ohne Kästchen (Diagramm zuordnen); Aussage mit „wahr / falsch“ zum Ankreuzen
\newcommand{\gbildname}[3]{\begin{minipage}[t][#3][b]{\gbildbreite}\centering #2\par\smallskip\mbox{\textbf{#1}}\end{minipage}}
\newcommand{\gwahrfalsch}[2]{\par\noindent\hangindent2.2em\hangafter1\makebox[2.2em][l]{(#1)}#2\hfill\mbox{wahr $\square$\quad falsch $\square$}\par}
\begin{document}

\pfheftstil
\pflaufkopf[P10 \pkt{} Prüfung]{Anteil und Veränderung in Prozent}{5ZQ1}\pfseitevon
\pfheftkopf{Anteil und Veränderung in Prozent}{}
\pfstufenzeile{Prozent und Anteil umwandeln \pkt{} Anteil in Prozent gesucht \pkt{} Veränderung in Prozent gesucht \pkt{} Aussagen prüfen \pkt{} Flächenanteil in Prozent gesucht}

% ===================== Stufe: Prozent und Anteil umwandeln =====================
\gkopfvor{\pfstufekopf{Prozent und Anteil umwandeln}{\pfjahre{22}{20}}}
% kein Einstieg: keine Schülerfrage in der Gliederung (Schülerfrage, -Fälle, -Form, -Weg, -Skizze; oder Frage) – Einstieg nur von dort, nie aus der Bank (pruefung.md § 5, 12.10.)
% 2022-OS-B1f-m1 (Original 2022-OS-B1f, Schwere I, rang 1; herausgelöste Fassung: keine)
\begin{pfaufg}{’22~\pkt{}~1f}{1.}{Basis}
Ein Zehntel aller Jugendlichen kauft nie ein Buch. Kreuze den passenden Prozentsatz an:
\par\smallskip\noindent\pfkreuz{A}{1 \%}\pfkreuz{B}{10 \%}\pfkreuz{C}{20 \%}\pfkreuz{D}{50 \%.}
\fusshilfe{\mbox{1:~B}}
\end{pfaufg}

% ===================== Stufe: Prozentsatz =====================
\Needspace*{0.16\textheight}\pfstufekopf{Anteil in Prozent gesucht}{\pfjahre{25 \pkt{} 23}{18 \pkt{} 15}}
% Schülerfrage der Stufe (Gliederung; Einstieg nur von dort, nie aus der Bank – 12.10.)
\begin{pfaufg}{}{2.}{}
In einer Klasse sind 20 Kinder, 5 davon kommen mit dem Rad. Was ist der Teil, was das Ganze? Wie viel Prozent kommen mit dem Rad?
\par\smallskip
\par\noindent Rad-Kinder in der Klasse\hfill\hfill Teil: \linie{30mm}\par\noindent\hfill Ganzes: \linie{30mm}\par\noindent\hfill Anteil: \linie{30mm}\,\%
\par\fusshilfe{\mbox{2:~Teil 5, Ganzes 20, 25 \%}}
\end{pfaufg}
% 2025-OS-B1e-m1 (Original 2025-OS-B1e, Schwere I, rang 11; herausgelöste Fassung: keine)
\begin{pfaufg}{’25~\pkt{}~1e}{3.}{Basis}
Bestimme, wie viel Prozent der Kreisfläche grau gefärbt sind.
\gzwei{\kreissektor[0.85]{144}{144°}}{\karo{3}}
\fusshilfe{\mbox{3:~40 \%}}
\end{pfaufg}
% 2023-OS-K6a-m1 (Original 2023-OS-K6a, Schwere I, rang 12; herausgelöste Fassung: keine)
\begin{pfaufg}{nach~P10\\’23~\pkt{}~6a}{4.}{}
Im Jahr 2022 gab es auf der Erde rund 8,0 Milliarden Menschen. Die Tabelle teilt sie nach dem Alter auf. Berechne den Anteil der Jugendlichen an der Weltbevölkerung im Jahr 2022 in Prozent.
\gzwei{\begin{tabular}{@{}lr@{}}\hline Altersgruppe & Anzahl in Mrd.\\\hline Kinder (0 bis 14 J.) & 2{,}0\\ Jugendliche (15 bis 24 J.) & 1{,}2\\ mittleres Alter (25 bis 64 J.) & 4{,}0\\ höheres Alter (ab 65 J.) & 0{,}8\\\hline\end{tabular}}{\vspace{2mm}\linie{30mm} \%}
\fusshilfe{\mbox{4:~15 \%}}
\end{pfaufg}

% ===================== Stufe: Erhöhung und Veränderung in Prozent =====================
\Needspace*{0.16\textheight}\pfstufekopf{Veränderung in Prozent gesucht}{\pfjahre{26 \pkt{} 22}{21 \pkt{} 20 \pkt{} 16}}
% Schülerfrage der Stufe (Gliederung; Einstieg nur von dort, nie aus der Bank – 12.10.)
\begin{pfaufg}{}{5.}{}
Ein Eis kostet erst 2 €, dann 2,50 €. Um wie viel Euro steigt der Preis? Welcher Preis ist das Ganze? Um wie viel Prozent steigt er?
\par\smallskip
\par\noindent Eis\hfill\hfill Unterschied: \linie{30mm}\,€\par\noindent\hfill Ganzes: \linie{30mm}\,€\par\noindent\hfill Steigerung: \linie{30mm}\,\%
\par\fusshilfe{\mbox{5:~0,50 €, alter Preis 2 €, 25 \%}}
\end{pfaufg}
% 2022-OS-K4b-m1 (Original 2022-OS-K4b, Schwere I, rang 7; herausgelöste Fassung: keine)
\begin{pfaufg}{nach~P10\\’22~\pkt{}~4b}{6.}{}
Das Säulendiagramm zeigt, wie viele von jeweils 1500 befragten Jugendlichen in den Jahren 2017 bis 2021 mehrmals pro Woche ein Buch lesen. Ermittle, um wie viel Prozent diese Zahl von 2020 auf 2021 zurückgegangen ist.
\gzwei{\begin{tikzpicture}[x=1.0cm,y=0.022cm,font=\scriptsize]\draw[->] (0,490)--(5.6,490);\draw[->] (0,490)--(0,615) node[above]{Anzahl};\draw (-0.08,500)--(0,500) node[left]{500};\draw[black!15] (0,500)--(5.4,500);\draw (-0.08,520)--(0,520) node[left]{520};\draw[black!15] (0,520)--(5.4,520);\draw (-0.08,540)--(0,540) node[left]{540};\draw[black!15] (0,540)--(5.4,540);\draw (-0.08,560)--(0,560) node[left]{560};\draw[black!15] (0,560)--(5.4,560);\draw (-0.08,580)--(0,580) node[left]{580};\draw[black!15] (0,580)--(5.4,580);\draw (-0.08,600)--(0,600) node[left]{600};\draw[black!15] (0,600)--(5.4,600);\fill[black!35] (0.55,490) rectangle (1.15,560);\node[above] at (0.85,560){560};\node[below] at (0.85,490){2017};\fill[black!35] (1.55,490) rectangle (2.15,530);\node[above] at (1.85,530){530};\node[below] at (1.85,490){2018};\fill[black!35] (2.55,490) rectangle (3.15,575);\node[above] at (2.85,575){575};\node[below] at (2.85,490){2019};\fill[black!35] (3.55,490) rectangle (4.15,585);\node[above] at (3.85,585){585};\node[below] at (3.85,490){2020};\fill[black!35] (4.55,490) rectangle (5.15,520);\node[above] at (4.85,520){520};\node[below] at (4.85,490){2021};\end{tikzpicture}}{\vspace{2mm}\linie{30mm} \%}
\fusshilfe{\mbox{6:~11,1 \%}}
\end{pfaufg}
% 2026-FOR-K3c-m1 (Original 2026-FOR-K3c, Schwere II, rang 8; herausgelöste Fassung: keine)
\begin{pfaufg}{nach~P10\\’26~\pkt{}~3c}{7.}{}
An einer Tankstelle wurde eine Woche lang täglich um 18 Uhr der Benzinpreis notiert (€ pro Liter): Mo 1,81, Di 1,82, Mi 1,80, Do 1,83, Fr 1,89, Sa 1,86, So 1,90. Berechne die prozentuale Preissteigerung von Donnerstag auf Freitag.
\hfill\linie{30mm} \%
\fusshilfe{\mbox{7:~3,3 \%}}
\end{pfaufg}

% ===================== Stufe: Aussagen prüfen =====================
\gkopfvor{\pfstufekopf{Aussagen prüfen}{\pfjahre{25 \pkt{} 23}{19 \pkt{} 17}}}
% kein Einstieg: keine Schülerfrage in der Gliederung (Schülerfrage, -Fälle, -Form, -Weg, -Skizze; oder Frage) – Einstieg nur von dort, nie aus der Bank (pruefung.md § 5, 12.10.)
% 2025-OS-K4b-m1 (Original 2025-OS-K4b, Schwere II, rang 9; herausgelöste Fassung: keine)
\begin{pfaufg}{nach~P10\\’25~\pkt{}~4b}{8.}{}
Eine Rampe für Rollstühle darf höchstens 8 \% Steigung haben, damit man sie ohne fremde Hilfe hinauffahren kann. Vervollständige den Satz: „Bei 8 \% Steigung steigt man auf \linie{25mm} Metern waagerechter Strecke um \linie{25mm} Meter an.“ Am Hauseingang soll eine Rampe eine 14 cm hohe Stufe überwinden; waagerecht ist sie 130 cm lang. Frau Yücel meint: „Unsere Rampe wird zu steil, ihre Steigung liegt über 8 \%.“ Überprüfe, ob Frau Yücel recht hat.
\gzwei{\begin{tikzpicture}[x=0.035cm,y=0.035cm,line width=0.8pt,font=\small]\draw (0,0)--(130,0)--(130,14)--cycle;\draw[line width=0.4pt] (124,0)--(124,6)--(130,6);\node[below] at (65,0){130 cm};\node[right] at (130,7){14 cm};\node[above left] at (65,7){Rampe};\end{tikzpicture}}{\karo{4}}
\fusshilfe{\mbox{8:~100 m und 8 m; ja}}
\end{pfaufg}
% 2023-OS-K6b-m1 (Original 2023-OS-K6b, Schwere II, rang 10; herausgelöste Fassung: keine)
\begin{pfaufg}{nach~P10\\’23~\pkt{}~6b}{9.}{}
Im Jahr 2022 gab es auf der Erde rund 8,0 Milliarden Menschen. Die Tabelle teilt sie nach dem Alter auf. Eine der beiden folgenden Aussagen widerspricht der Tabelle: „2022 war jeder zweite Mensch auf der Erde im mittleren Alter.“ „2022 war ein Drittel aller Menschen im Kindesalter.“ Kreuze die falsche Aussage an. Formuliere eine berichtigte Aussage, die zur Tabelle passt.
\gzwei{\begin{tabular}{@{}lr@{}}\hline Altersgruppe & Anzahl in Mrd.\\\hline Kinder (0 bis 14 J.) & 2{,}0\\ Jugendliche (15 bis 24 J.) & 1{,}2\\ mittleres Alter (25 bis 64 J.) & 4{,}0\\ höheres Alter (ab 65 J.) & 0{,}8\\\hline\end{tabular}}{\karo{3}}
\fusshilfe{\mbox{9:~falsch; 0,25, ein Viertel aller Menschen}}
\end{pfaufg}

% ===================== Stufe: Prozent aus einer berechneten Fläche =====================
\gkopfvor{\pfstufekopf{Flächenanteil in Prozent gesucht}{\pfjahre{23}{}}}
% kein Einstieg: keine Schülerfrage in der Gliederung (Schülerfrage, -Fälle, -Form, -Weg, -Skizze; oder Frage) – Einstieg nur von dort, nie aus der Bank (pruefung.md § 5, 12.10.)
% 2023-OS-K5c-m1 (Original 2023-OS-K5c, Schwere III, rang 11; herausgelöste Fassung: keine)
\begin{pfaufg}{nach~P10\\’23~\pkt{}~5c}{10.}{FOR}
Für Konservendosen werden kreisrunde Deckel mit dem Radius 3,5 cm aus Blechstreifen ausgestanzt, die 7 cm breit und 1 m lang sind. Aus jedem Streifen sollen möglichst viele Deckel entstehen. Der Abfall darf höchstens 24 \% eines Streifens ausmachen. Überprüfe rechnerisch, ob diese Vorgabe eingehalten wird.
\gzwei{\begin{tikzpicture}[line width=0.8pt,font=\small]\draw (0,0) rectangle (5.0,0.5);\foreach \x in {0.25,0.75,1.25,1.75} {\draw[line width=0.5pt] (\x,0.25) circle (0.25);}\node at (2.4,0.25) {$\cdots$};\draw[line width=0.5pt] (4.75,0.25) circle (0.25);\node[below] at (2.5,0) {1 m};\node[right] at (5.0,0.25) {7 cm};\end{tikzpicture}}{\karo{4}}
\fusshilfe{\mbox{10:~14 Deckel; 161,2 cm²; ja}}
\end{pfaufg}

\pfblattende{11}{26 FOR & 25 & 24 & 23 & 22 & 21 & 20 & 19 & 18 & 17 & 16}{3c & 1e 4b &  & 5c 6a 6b & 1f 4b & 5a & 1a 2d 4c & 5b & 7a & 2b & 2c}%
{\pfausschnitt{}{Anteil und Veränderung in Prozent}{Prozentwert und Grundwert}}
\end{document}